\chapter{Resultados y discusión}

% ============================================================
\section{Condiciones de simulación y casos de estudio}
% ============================================================

\subsection{Condiciones generales de simulación}

\subsection{Escenarios de perturbación}

Se consideran los siguientes escenarios:

\begin{enumerate}
    \item Sin torque de carga y sin cogging.
    \item Con torque de carga y sin cogging.
    \item Sin torque de carga y con cogging.
    \item Con torque de carga y con cogging.
\end{enumerate}

\subsection{Configuraciones del observador}

Se comparan las siguientes configuraciones:

\begin{enumerate}
    \item Observador base de cuatro estados.
    \item Observador de cinco estados con torque equivalente.
    \item Observador de estados aumentados con componentes armónicas.
\end{enumerate}


% ============================================================
\section{Validación del sistema de control}
% ============================================================

\subsection{Respuesta de los lazos de corriente}

\subsection{Respuesta del lazo de velocidad}

\subsection{Efecto del torque de carga y del cogging}


% ============================================================
\section{Desempeño del observador base de cuatro estados}
% ============================================================

\subsection{Operación sin perturbaciones}

\subsection{Respuesta frente a torque de carga}

\subsection{Respuesta frente al cogging}

\subsection{Respuesta frente a carga y cogging simultáneos}

\subsection{Error de estimación de posición y velocidad}


% ============================================================
\section{Desempeño del observador de cinco estados}
% ============================================================

\subsection{Estimación del torque equivalente}

\subsection{Respuesta frente a torque de carga}

\subsection{Respuesta frente al cogging}

\subsection{Respuesta frente a carga y cogging simultáneos}

\subsection{Efecto sobre la estimación de posición y velocidad}


% ============================================================
\section{Desempeño del observador de estados aumentados}
% ============================================================

\subsection{Incorporación de componentes armónicas}

\subsection{Reconstrucción de las perturbaciones periódicas}

\subsection{Respuesta frente a torque de carga}

\subsection{Respuesta frente al cogging}

\subsection{Respuesta frente a carga y cogging simultáneos}

\subsection{Efecto del número de estados aumentados}


% ============================================================
\section{Comparación de las configuraciones del observador}
% ============================================================

\subsection{Error de estimación de posición}

\subsection{Error de estimación de velocidad}

\subsection{Reconstrucción del torque perturbador}

\subsection{Atenuación y desfase de la velocidad estimada}

\subsection{Síntesis comparativa}


% ============================================================
\section{Operación sensorless a baja velocidad}
% ============================================================

\subsection{Respuesta de la inyección de alta frecuencia}

\subsection{Estimación sensorless de posición}

\subsection{Estimación sensorless de velocidad}

\subsection{Observador base en operación sensorless}

\subsection{Observador de cinco estados en operación sensorless}

\subsection{Observador de estados aumentados en operación sensorless}

\subsection{Efecto del torque de carga y del cogging}


% ============================================================
\section{Análisis de observabilidad}
% ============================================================

\subsection{Observabilidad del modelo base}

\subsection{Efecto de la velocidad y la saliencia}

\subsection{Efecto de la inyección de alta frecuencia}

\subsection{Observabilidad de los modelos aumentados}


% ============================================================
\section{Discusión de resultados}
% ============================================================

\subsection{Reconstrucción de los estados mecánicos}

\subsection{Reconstrucción de las perturbaciones}

\subsection{Compromiso entre complejidad y capacidad de estimación}

\subsection{Limitaciones de la operación sensorless}